\documentclass[11pt]{article}

\usepackage[margin=1in]{geometry}
\usepackage{booktabs}
\usepackage{amsmath}
\usepackage{amssymb}
\usepackage{graphicx}
\usepackage{hyperref}
\usepackage{xcolor}

\newcommand{\todo}[1]{\textcolor{red}{[TODO: #1]}}
\newcommand{\method}{Base+Reasoner CP}

\title{Reasoning-Step Verification Features for Conformal Legal Issue Prediction}
\author{Anonymous}
\date{}

\begin{document}
\maketitle

\begin{abstract}
Legal issue prediction systems are often evaluated as fixed-label classifiers, but in legal practice a reliable set of plausible issues can be more useful than a single brittle prediction. We study conformal multi-label prediction for ECtHR Article prediction, where the goal is to output a calibrated set of Convention articles that contains all gold articles with a target probability. Our central question is whether noisy reasoning-step verification signals from a LegalReasoner-style pipeline can improve the efficiency of conformal prediction, without treating those reasoning steps as ground-truth explanations.

We propose a two-stage scoring approach. First, a base legal classifier produces per-label probabilities. Second, for the top candidate labels, a structured LLM pipeline generates short article-specific reasoning steps and verifies whether each step is supported by the case facts. The resulting verifier features are merged with base classifier scores and used to train a supervised per-label meta-scorer. The meta-scorer is then calibrated using split conformal prediction.

In a 100/100/100 ECtHR-B pilot, the combined conformal predictor matched the coverage of a classifier-only conformal predictor at $\alpha=0.1$ while returning smaller prediction sets: average set size decreased from 2.39 to 2.27, with micro-F1 increasing from 0.668 to 0.696. These pilot results are preliminary, but they support the paper's core hypothesis: reasoning-step verification signals, even when noisy and not trusted as proof, can provide calibrated evidence that improves legal prediction-set efficiency. Full-scale experiments and stronger encoder baselines remain necessary before making a final empirical claim.
\end{abstract}

\section{Introduction}

Legal issue prediction is typically framed as multi-label classification: given the facts of a case, predict which legal provisions are implicated. In the European Court of Human Rights (ECtHR) setting, this means predicting one or more Convention articles such as Article 3, Article 6, Article 8, or Protocol 1 Article 1. Standard classifiers return point predictions or thresholded label sets, but legal use cases often need a different guarantee: if a system returns a set of plausible legal issues, how likely is it that the set contains all issues annotated by the benchmark?

Conformal prediction provides a natural language for this problem. Instead of claiming that a model's top label is correct, a conformal predictor returns a set $C_\alpha(x)$ such that, under standard exchangeability assumptions,
\[
  \Pr[Y \subseteq C_\alpha(X)] \geq 1-\alpha,
\]
where $Y$ is the set of gold labels. This set-valued output is especially attractive in law, where missing a relevant issue can be costly and where calibrated uncertainty is more useful than overconfident point prediction.

This paper asks whether reasoning-style LLM signals can make those conformal sets smaller or more useful at the same target coverage. Recent LLM legal reasoning pipelines can generate candidate issues, reason over them, and verify the reasoning against the input facts. However, treating generated reasoning steps as ground-truth explanations is risky: the model may hallucinate, overgeneralize, or verify a plausible but legally irrelevant chain. We therefore make a more modest claim. We do not claim that the LLM proves each reasoning step true. Instead, we ask whether noisy reasoning and verification signals are useful features for calibrated legal issue prediction.

Our method combines a base legal classifier with LegalReasoner-style verifier features. The base classifier scores every label. The LLM is only asked to reason over the top-$K$ candidate labels from the base classifier. For labels not sent to the LLM, verifier features are set to defaults, so any label can still enter the final conformal set through the base classifier score. We train three meta-scorers: a base-only scorer, a reasoner-only scorer, and a combined base-plus-reasoner scorer. All are evaluated using the same multi-label conformal calibration procedure.

The intended paper-grade claim is:

\begin{quote}
At fixed target coverage, a conformal predictor using reasoning-step verification features returns smaller or more useful legal article sets than classifier-only, LLM-only, and fixed-threshold baselines.
\end{quote}

The present draft reports a small pilot and lays out the full experimental design. The pilot result is promising but not final: on 100 ECtHR-B test cases, the combined predictor matches classifier-only coverage while reducing average conformal set size from 2.39 to 2.27. The remaining work is to scale the experiment, add strong encoder baselines, and conduct ablations showing that verifier features add signal beyond base-score rank.

\section{Task and Data}

We use the ECtHR-B configuration of LexGLUE. Each example consists of case facts and a multi-label set of Convention articles. The standard splits are:
\begin{itemize}
  \item train: supervised training for the base classifier and meta-scorer;
  \item validation: conformal calibration;
  \item test: final held-out evaluation.
\end{itemize}

Let $\mathcal{Y}$ be the set of legal article labels. In the current ECtHR-B setup used by the repository, the label set contains:
\[
\{2,3,5,6,8,9,10,11,14,\mathrm{P1\mbox{-}1}\}.
\]
For each case $x_i$, the gold label set is $Y_i \subseteq \mathcal{Y}$.

\section{Method}

\subsection{Base Classifier}

The base scorer is a reproducible text classifier:
\begin{itemize}
  \item input: normalized case text built from the case paragraphs;
  \item representation: TF-IDF unigrams and bigrams;
  \item classifier: one-vs-rest logistic regression;
  \item output: one probability per legal label.
\end{itemize}

For each case and label, the base classifier returns $p_{\mathrm{base}}(y \mid x)$. These scores serve two purposes. First, they form a classifier-only conformal baseline. Second, the top-$K$ labels are passed to the LegalReasoner-style feature generator.

\subsection{Reasoning and Verification Features}

For each case, we select the top-$K$ labels by base classifier probability. In the pilot, $K=6$. For each selected candidate label, a structured LLM pipeline performs two stages:

\paragraph{Reasoning stage.}
The model writes short fact-grounded reasoning steps for each candidate label and assigns a preliminary decision:
\[
\{\mathrm{support}, \mathrm{weak}, \mathrm{none}\}.
\]

\paragraph{Verification stage.}
The model checks whether each reasoning step is supported by the case facts and assigns:
\[
\{\mathrm{supported}, \mathrm{partial}, \mathrm{unsupported}\}
\]
both at the step level and at the label level.

The resulting record is flattened into per-label numeric features, including:
\begin{itemize}
  \item whether the label was in the top-$K$ candidate set;
  \item normalized candidate rank;
  \item preliminary support/weak/none indicators;
  \item verifier supported/partial/unsupported indicators;
  \item number of reasoning steps;
  \item number of supported, partial, and unsupported step checks;
  \item supported-step fraction;
  \item evidence string length;
  \item hard support gate indicator.
\end{itemize}

Labels outside the top-$K$ candidate set receive default verifier features. This is important: the LLM is not allowed to remove labels from consideration. Any label can still enter the final conformal set through the base classifier probability.

\subsection{Meta-Scorers}

We train per-label supervised meta-scorers on rows of the form $(x_i,y)$, with target
\[
  t_{i,y} = \mathbb{1}[y \in Y_i].
\]
The current implementation uses logistic regression over dictionary features. We train three scorers:
\begin{enumerate}
  \item \textbf{Base-only}: uses only $p_{\mathrm{base}}(y \mid x)$ and label identity;
  \item \textbf{Reasoner-only}: uses verifier/reasoning features and label identity, excluding the base probability;
  \item \textbf{Base+Reasoner}: uses both base probability and verifier/reasoning features.
\end{enumerate}

The output score is a calibrated probability-like score
\[
  p_{\mathrm{meta}}(y \mid x).
\]

\section{Conformal Calibration}

We use split conformal prediction on the validation split. For multi-label set coverage, define the calibration nonconformity score for case $i$ as:
\[
  S_i = \max_{y \in Y_i} \left(1 - p_i(y)\right),
\]
where $p_i(y)$ is the score assigned to label $y$ for case $i$.

Let $q_\alpha$ be the conformal quantile of the calibration scores:
\[
  q_\alpha = \mathrm{Quantile}_{\lceil(n+1)(1-\alpha)\rceil/n}(\{S_i\}_{i=1}^n).
\]
The prediction set is:
\[
  C_\alpha(x) = \{y \in \mathcal{Y}: 1 - p(y \mid x) \leq q_\alpha \}.
\]

This gives finite-sample marginal coverage under exchangeability:
\[
  \Pr[Y_{\mathrm{test}} \subseteq C_\alpha(X_{\mathrm{test}})] \geq 1-\alpha.
\]

We also implement labelwise, or Mondrian-style, thresholds:
\[
  q_{\alpha,y} = \mathrm{Quantile}(\{1-p_i(y): y \in Y_i\}),
\]
with a fallback global threshold for labels with insufficient calibration examples. Labelwise calibration is especially relevant in legal prediction because rare articles can matter disproportionately.

\section{Experimental Setup}

\subsection{Current Pilot}

The current pilot uses:
\begin{itemize}
  \item 100 training cases for the meta-scorer;
  \item 100 validation cases for conformal calibration;
  \item 100 test cases for evaluation;
  \item top-$K=6$ base classifier candidates for LLM reasoning;
  \item Qwen2.5-7B-Instruct served through vLLM;
  \item TF-IDF one-vs-rest logistic regression as the base classifier.
\end{itemize}

The reasoner feature generation is run on a GPU cluster through Slurm. Rows with malformed or incomplete LLM outputs are repaired by rerunning only failed case IDs and merging successful repairs into the split JSONL file.

\subsection{Planned Full-Scale Setting}

The paper-final experiment should use:
\begin{itemize}
  \item full validation split for calibration;
  \item full test split for final evaluation;
  \item a large or full train split for meta-scorer training;
  \item sharded reasoner feature generation for the train split;
  \item a documented repair policy for malformed structured LLM outputs.
\end{itemize}

The full-scale experiment should also include a stronger base encoder, such as Legal-BERT or Longformer, to ensure that improvements are not limited to a weak TF-IDF classifier. \todo{Add transformer base scorer and compare with TF-IDF baseline.}

\subsection{Evaluation Metrics}

For each method and each $\alpha$, we report:
\begin{itemize}
  \item set coverage: fraction of cases where $Y_i \subseteq C_\alpha(x_i)$;
  \item average prediction-set size;
  \item median prediction-set size;
  \item micro-F1 and macro-F1 when treating the conformal set as a prediction;
  \item micro and macro recall;
  \item per-label recall;
  \item rare-label recall;
  \item coverage by number of gold labels;
  \item coverage by case-length bucket.
\end{itemize}

The first two metrics, coverage and average set size, define the main conformal tradeoff.

\section{Pilot Results}

Table~\ref{tab:pilot-global} reports the current 100/100/100 pilot at $\alpha=0.1$ with global conformal thresholds.

\begin{table}[h]
\centering
\begin{tabular}{lccccc}
\toprule
Method & Coverage & Avg. $|C|$ & Median $|C|$ & Micro-F1 & Macro-F1 \\
\midrule
Base CP & 0.910 & 2.39 & 2.00 & 0.668 & 0.682 \\
Reasoner-only CP & 0.870 & 2.56 & 2.00 & 0.619 & 0.608 \\
Base+Reasoner CP & 0.910 & 2.27 & 2.00 & 0.696 & 0.706 \\
\bottomrule
\end{tabular}
\caption{Current ECtHR-B pilot results at $\alpha=0.1$ using global conformal calibration. This is a 100 train / 100 validation / 100 test pilot and should not be treated as a final result.}
\label{tab:pilot-global}
\end{table}

The combined model matches the coverage of the base-only conformal predictor while reducing average set size from 2.39 to 2.27. It also improves micro-F1 and macro-F1. The reasoner-only model performs worse, which suggests that the verifier signal is not sufficient by itself but can be useful when combined with the base classifier.

This is the desired qualitative pattern. However, the effect size is modest and the sample is small. A 0.12-label reduction in average set size is promising but not enough for a strong paper without full-scale replication, confidence intervals, and ablations.

\section{Baselines and Ablations}

The full paper should compare against:
\begin{enumerate}
  \item base classifier top-$k$;
  \item base classifier fixed threshold;
  \item base classifier conformal prediction;
  \item direct LLM label prediction;
  \item LegalReasoner hard gate;
  \item reasoner-only conformal prediction;
  \item base-plus-reasoner conformal prediction.
\end{enumerate}

The most important ablations are:
\begin{itemize}
  \item base probability only;
  \item base probability plus candidate rank only;
  \item base probability plus preliminary decision;
  \item base probability plus verifier verdict;
  \item base probability plus step-count features;
  \item base probability plus full verifier features.
\end{itemize}

These ablations address the main reviewer concern: the combined model may improve only because it receives candidate-rank information derived from the base classifier, not because reasoning-step verification adds signal. The central paper claim requires showing that verifier features improve over rank-only features.

\section{Discussion}

The pilot supports a modest but meaningful hypothesis. LLM reasoning steps are not treated as explanations that prove legal truth. Instead, they are treated as noisy structured evidence. Conformal calibration then absorbs uncertainty in the score function and produces sets with target coverage.

This framing has three advantages:
\begin{itemize}
  \item It avoids overclaiming the faithfulness of generated reasoning.
  \item It uses the LLM only as a feature generator, not as an uncalibrated decision maker.
  \item It measures success through coverage-efficiency tradeoffs rather than raw point-prediction accuracy.
\end{itemize}

The pilot also shows why a hard LLM gate is dangerous. The support-only gate often misses gold labels. In the current pilot, the hard support signal is useful but too sparse to serve as the final prediction rule. The combined conformal model can use the verifier when it helps, while still allowing the base classifier to include labels that the LLM did not support.

\section{Limitations}

This draft has several limitations:
\begin{itemize}
  \item The current results are based on a small pilot and are not paper-final.
  \item The base classifier is TF-IDF logistic regression, not a strong transformer legal encoder.
  \item Structured LLM generation occasionally fails and requires a repair policy.
  \item Rare labels have too few pilot examples for stable conclusions.
  \item Ablations are not yet implemented.
  \item The current conformal evaluation is marginal; labelwise calibration needs more careful treatment for rare labels.
\end{itemize}

\section{Conclusion}

We propose using reasoning-step verification signals as features for conformal legal issue prediction. The goal is not to prove each generated step true, but to test whether noisy verifier signals can improve calibrated multi-label prediction sets. A small ECtHR-B pilot suggests that combining base classifier probabilities with LegalReasoner-style verifier features can match classifier-only coverage while reducing average set size. The result is promising but preliminary. Full-scale evaluation, strong encoder baselines, rare-label analysis, and feature ablations are necessary before making a final publication claim.

\begin{thebibliography}{9}

\bibitem{vovk2005}
Vladimir Vovk, Alex Gammerman, and Glenn Shafer.
\newblock \emph{Algorithmic Learning in a Random World}.
\newblock Springer, 2005.

\bibitem{angelopoulos2021}
Anastasios N. Angelopoulos and Stephen Bates.
\newblock A gentle introduction to conformal prediction and distribution-free uncertainty quantification.
\newblock arXiv preprint, 2021.

\bibitem{lexglue}
Ilias Chalkidis, Abhik Jana, Dirk Hartung, Michael Bommarito, Ion Androutsopoulos, Daniel Katz, and Nikolaos Aletras.
\newblock LexGLUE: A benchmark dataset for legal language understanding in English.
\newblock In \emph{Proceedings of ACL}, 2022.

\bibitem{legalbert}
Ilias Chalkidis, Manos Fergadiotis, Prodromos Malakasiotis, Nikolaos Aletras, and Ion Androutsopoulos.
\newblock LEGAL-BERT: The muppets straight out of law school.
\newblock In \emph{Findings of EMNLP}, 2020.

\end{thebibliography}

\end{document}
